\subsection{AI-Assisted Learning: The 2025--2026 Context}

The rise of AI development tools has become a practical factor shaping how independent builders learn. Tools such as GitHub Copilot and general-purpose assistants can compress the time between a question and a plausible next step by summarizing documentation, generating draft code, and explaining unfamiliar concepts.

\subsubsection{Documentation Acceleration}

AI tools can reduce search and synthesis time by providing immediate summaries of APIs, error messages, and implementation patterns. This often increases momentum during project work because the learner can iterate quickly on a hypothesis. However, perceived speed is not always equivalent to actual productivity on complex tasks, and AI-generated suggestions can require non-trivial verification and rework.

\subsubsection{Validation Through AI Code Review}

AI tools can function as a lightweight review partner by flagging potential bugs, suggesting refactors, and proposing tests. For a solo learner without continuous access to expert feedback, this can partially approximate peer review. Nevertheless, AI review remains heuristic: it can miss deeper architectural issues and may offer changes that appear idiomatic but are misaligned with the project constraints.

\subsubsection{Responsible AI Usage as a Skill}

Effective use of AI systems is itself a learned capability. Learners must decide when to consult AI, how to specify constraints, and how to evaluate outputs. In practice, this shifts part of technical competence toward prompt formulation, verification habits, and the ability to detect confident but incorrect outputs.

\subsubsection{Risks and Failure Modes: Hallucinations, Dependency, and Cognitive Outsourcing}

AI-assisted learning also introduces risks that can reduce learning quality and threaten academic credibility if left unaddressed.

\begin{enumerate}
\item \textbf{Hallucinations and fabricated sources.} Language models can generate confident explanations, claims, or citations that are incorrect or do not exist. In academic writing, this is especially damaging because a single fabricated citation can undermine the credibility of the overall work. This implies that AI outputs should be treated as drafts rather than evidence and that claims must be grounded in primary sources.
\item \textbf{Dependency and workflow fragility.} Heavy reliance on AI can make learning workflows brittle. Changes in access (pricing, policy, outages) or model behavior after updates can reduce reproducibility and limit the learner's ability to work independently.
\item \textbf{Shallow understanding and the illusion of competence.} AI can produce working solutions without requiring the learner to form an internal model of why the solution works. This may reduce transfer to unfamiliar problems and can conceal gaps until a novel failure occurs.
\item \textbf{Cognitive outsourcing and skill atrophy.} Offloading planning, recall, and explanation to AI can weaken core practices associated with expertise development, such as deliberate practice, self-explanation, and reflective debugging \citep{ericsson1993}. Used as scaffolding, AI can support learning; used as substitution, it can erode independent problem solving.
\end{enumerate}

These risks do not negate the utility of AI tools; rather, they imply that AI-assisted learning should be treated as a bounded technique with explicit guardrails, especially in research contexts where credibility and verifiability are central \citep{bender2021}.
